\documentclass[12pt, a4paper]{article}

% Encodage et langue
\usepackage[utf8]{inputenc}
\usepackage[T1]{fontenc}
\usepackage[french]{babel}

% Mise en page standard
\usepackage{geometry}
\geometry{hmargin=2.5cm, vmargin=2.5cm}
\usepackage{xcolor}

% Remplacement de 'parskip' par des commandes standard
\setlength{\parindent}{0pt}
\setlength{\parskip}{1em}

% Définition de la couleur (Bleu)
\definecolor{slideheader}{RGB}{30, 144, 255}

% Commande simplifiée pour les en-têtes (sans le paquet calc)
\newcommand{\slideheader}[2]{
    \vspace{1cm}
    \noindent
    \colorbox{slideheader!20}{
        % Utilisation de dimexpr natif pour éviter le paquet calc
        \parbox{\dimexpr\linewidth-2\fboxsep\relax}{
            \textbf{\Large \textcolor{black}{Diapositive #1 : #2}}
        }
    }
    \vspace{0.5cm}
}

% Informations du document
\title{Script de Soutenance Oral\\ \large Gestion de projet informatique}
\author{Arnaud Gadille, Maël Mignard}
\date{}

\begin{document}

\maketitle

\section*{Introduction}
\textit{(Transition après l'introduction générale du projet)}

\slideheader{3}{Bases, acteurs et outils}

«~Pour commencer à structurer notre approche, penchons-nous sur les fondations de la gestion de projet. Sur cette diapositive, nous avons synthétisé les trois piliers essentiels.

D'abord, les \textbf{Bases}.\\
Gérer un projet, ce n'est pas improviser. C'est un processus cyclique qui consiste à planifier, exécuter, suivre et enfin clôturer une action pour atteindre un but précis.

Ensuite, les \textbf{Acteurs}.\\
Un projet ne se fait jamais seul. On distingue trois entités clés :
\begin{itemize}
    \item Le \textbf{Chef de projet}, qui pilote le navire.
    \item L'\textbf{Équipe technique}, qui réalise la production.
    \item Et les \textbf{Clients}, pour qui le projet est réalisé.
\end{itemize}

Enfin, les \textbf{Outils}.\\
Pour faire le lien entre eux, on utilise des logiciels spécifiques (Jira, Trello, MS Project). Ils sont indispensables pour organiser les tâches et suivre l'avancement.~»

\slideheader{4}{Méthodes traditionnelles (Phases)}

«~Entrons dans le vif du sujet avec les \textbf{Méthodes Traditionnelles} (dites "Waterfall").
Comme vous le voyez, elles reposent sur une structure linéaire en \textbf{5 phases} :

\begin{enumerate}
    \item \textbf{L'Initiation} : Définir le projet et sa faisabilité.
    \item \textbf{La Planification} : Tout écrire (budget, calendrier) avant de commencer.
    \item \textbf{L'Exécution} : Produire en suivant le plan.
    \item \textbf{Le Suivi et Contrôle} : Surveiller les écarts de budget ou de temps.
    \item \textbf{La Clôture} : Livrer le produit final.
\end{enumerate}

L'objectif ici est la \textbf{prévisibilité} pour éviter les surprises.~»

\slideheader{5}{Autres méthodes (Cycle en V vs PERT)}

«~Comparons maintenant deux modèles spécifiques.

\medskip
\textit{(Geste vers la gauche)}
\textbf{Le Cycle en V} : C'est le standard pour les besoins stables. Chaque étape de conception (descendante) correspond à une phase de test (montante). C'est très carré, mais rigide : revenir en arrière coûte cher.

\medskip
\textit{(Geste vers la gauche)}
\textbf{Le PERT} : Ici, on voit le projet comme un \textbf{réseau de tâches}. On gère les dépendances (B ne commence pas avant A). Cela permet de trouver le \textbf{"chemin critique"}, c'est-à-dire la durée minimale incompressible du projet.~»

\end{document}